\documentclass[aspectratio=169]{beamer}
\usepackage[utf8]{inputenc}
\usepackage[spanish]{babel}
\usepackage{amsmath}
\usepackage{amssymb}
\usepackage{graphicx}
\usepackage{booktabs}
\usepackage{multirow}
\usepackage{xcolor}

% Tema y colores
\usetheme{Madrid}
\usecolortheme{default}
\setbeamertemplate{navigation symbols}{}
\setbeamertemplate{caption}[numbered]

% Definir colores personalizados
\definecolor{udea}{RGB}{0,102,51}
\setbeamercolor{structure}{fg=udea}
\setbeamercolor{title}{fg=white,bg=udea}
\setbeamercolor{frametitle}{fg=white,bg=udea}
\setbeamercolor{titlelike}{fg=white,bg=udea}
\setbeamercolor{author}{fg=black}
\setbeamercolor{institute}{fg=black}
\setbeamercolor{date}{fg=black}

% Información del documento
\title[Partículas en una Caja Caliente]{Miniproyecto 2: Where are the particles when the box is hot?}
\subtitle{Reproducción y Extensión al Pozo Finito}
\author[Juan Montoya \& Pablo Sánchez]{Juan Montoya \& Pablo Sánchez}
\institute[UdeA]{Instituto de Física\\Universidad de Antioquia}
\date{21 de Enero de 2026}

\begin{document}

%===============================================================================
% DIAPOSITIVA 1: PORTADA
%===============================================================================
\begin{frame}
\titlepage
\end{frame}

%===============================================================================
% DIAPOSITIVA 2: CONTENIDO
%===============================================================================
\begin{frame}{Contenido}
\tableofcontents
\end{frame}

%===============================================================================
% DIAPOSITIVA 3: INTRODUCCIÓN Y OBJETIVO
%===============================================================================
\section{Introducción}
\begin{frame}{Introducción y Objetivo}
\begin{block}{Objetivo Principal}
Reproducir los resultados del artículo \textit{"Where are the particles when the box is hot?"} (Eur. J. Phys., 2019) y extender el análisis al caso de un pozo de potencial finito.
\end{block}

\vspace{0.2cm}
\begin{columns}[T]
\column{0.5\textwidth}
\textbf{Problema central:}
\begin{itemize}
\item ¿Cómo cambia la densidad de probabilidad de una partícula cuántica al aumentar la temperatura?
\item ¿Qué efecto tiene la estadística (Bose-Einstein vs Fermi-Dirac) en sistemas de dos partículas?
\item ¿Cómo influye el espín en la distribución espacial?
\end{itemize}

\column{0.5\textwidth}
\textbf{Extensión propuesta:}
\begin{itemize}
\item Considerar un pozo de altura finita $V_0$.
\item Evaluar el límite termodinámico y la densidad de estados.
\end{itemize}
\end{columns}
\end{frame}

%===============================================================================
% SECCIÓN 1: POZO INFINITO (MODELO DEL ARTÍCULO)
%===============================================================================
\section{El Pozo de Potencial Infinito}

% Frame: Modelo de una partícula
\begin{frame}{1. Una partícula en el Pozo Infinito}
\begin{columns}[T]
\column{0.5\textwidth}
\vspace{-0.3cm}
\begin{block}{Modelo Básico}
Partícula de masa $m$ en una caja de longitud $L$:
$$ V(X) = \begin{cases} 0 & 0 < X < L \\ \infty & \text{resto} \end{cases} $$
\end{block}

\textbf{Soluciones:}
$$ \psi_n(X) = \sqrt{\frac{2}{L}} \sin\left(\frac{n\pi X}{L}\right) $$
$$ E_n = \frac{h^2 n^2}{8mL^2} $$

\column{0.5\textwidth}
\vspace{0.7cm}
\begin{center}
\includegraphics[width =\columnwidth]{figures/pozo_infinito_termico.eps}
\end{center}
\end{columns}
\end{frame}

\begin{frame}{1. Una partícula en el Pozo Infinito con temperatura}
\begin{columns}[T]
\column{0.3\textwidth}
\vspace{1cm}
\textbf{Ensamble Canónico (Temperatura $T$):}
Densidad de probabilidad térmica media ponderada por Boltzmann:
$$ P_{th}(x, t) = \frac{\sum_{n=1}^{\infty} |\psi_n(x)|^2 e^{-E_n/k_B T}}{\sum_{n=1}^{\infty} e^{-E_n/k_B T}}$$

\column{0.7\textwidth}
\vspace{-0.73cm}
\begin{center}
\includegraphics[width =0.55\columnwidth]{figures/fig1a_Pth.png}
\includegraphics[width =0.55\columnwidth]{figures/fig1b_boundary_layer.eps}

\end{center}
\end{columns}
\end{frame}


% Frame: Dos partículas (Bosones vs Fermiones)
\begin{frame}{2. Dos Partículas: Bosones vs Fermiones (Sin Espín)}
\begin{columns}[T]
\column{0.5\textwidth}
\vspace{-0.8cm}
\textbf{Bosones ($\Psi$ Simétrica):}
$$ \Psi = \frac{1}{2} ( \psi_{n_1}(x_1)\psi_{n_2}(x_2) + \psi_{n_1}(x_2)\psi_{n_2}(x_1) )$$
\begin{itemize}
\vspace{-0.8cm}
\item \textcolor{blue}{Atracción efectiva}: Pueden ocupar el mismo estado ($n_1 = n_2$).
\end{itemize}

\begin{center}
    \includegraphics[width = 1.04\columnwidth]{figures/figure2_bosons.png}
\end{center}

\column{0.5\textwidth}
\vspace{-0.8cm}
\textbf{Fermiones ($\Psi$ Antisimétrica):}
$$ \Psi = \frac{1}{2} ( \psi_{n_1}(x_1)\psi_{n_2}(x_2) - \psi_{n_1}(x_2)\psi_{n_2}(x_1) )$$
\begin{itemize}
\vspace{-0.8cm}
\item \textcolor{red}{Repulsión Pauli}: No pueden ocupar el mismo estado ($n_1 \neq n_2$).
\end{itemize}
\begin{center}
    \includegraphics[width = 1.04\columnwidth]{figures/figure3_fermions.png}
\end{center}
\end{columns}
\end{frame}

% Frame: Figura 4 - Densidad Térmica 2 Partículas
\begin{frame}{3. Densidad térmica de dos partículas}
\begin{columns}[T]
\column{0.55\textwidth}
\scriptsize
\begin{itemize}
    \item \textbf{Sumas Diferenciadas:}
    \begin{itemize}
        \item \textit{Bosones:} $1 \le n_1 \le n_2 \le n_{max}$ (Permite $n_1=n_2$).
        \item \textit{Fermiones:} $1 \le n_1 < n_2 \le n_{max}$ (Exclusion de Pauli).
    \end{itemize}
    \item \textbf{Convergencia:} $n_{max} \approx 25$ para estabilidad a $t=10$.
    \item
    \textbf{Bosones:} Máxima probabilidad en la diagonal $x_1=x_2$ (atracción efectiva).
    \item 
    \textbf{Fermiones (sin espín):} Probabilidad nula en diagonal.
    \item
    \textbf{Fermiones (con espín):} Al incluir espín, se permiten contribuciones simétricas espaciales (Singlete) que pueblan la diagonal.
\end{itemize}
\column{0.45\textwidth}
\vspace{-1.7cm}
\begin{center}
    \includegraphics[width=0.7\columnwidth]{figures/fig4_thermalized.eps}
\end{center}
\end{columns}
\end{frame}

% Frame: Fermiones con Espín
\begin{frame}{4. Mecanismo de Espín en Fermiones}
\vspace{-0.6cm}
\begin{columns}[T]
\column{0.45\textwidth}
\begin{block}{Singlete ($S=0$)}
\begin{itemize}
\item Espín Antisimétrico ($\uparrow\downarrow - \downarrow\uparrow$)
\item \textbf{Espacio Simétrico} (Tipo Bosón)
\end{itemize}
\end{block}

\column{0.5\textwidth}
\begin{block}{Triplete ($S=1$)}
\begin{itemize}
\item Espín Simétrico ($\uparrow\uparrow$, etc.)
\item \textbf{Espacio Antisimétrico} (Tipo Fermión)
\end{itemize}
\end{block}
\end{columns}

\vspace{-0.01cm}
\textbf{Resultado Térmico:}
$$ P_{total} = \frac{1}{4} P_{Singlete} + \frac{3}{4} P_{Triplete} $$

\vspace{-0.3cm}
\begin{center}
    \includegraphics[width=0.65\columnwidth]{figures/fig5_infinito.eps}
\end{center}

\end{frame}

%===============================================================================
% SECCIÓN 2: POZO FINITO (EXTENSIÓN)
%===============================================================================
\section{Extensión: Pozo de Potencial Finito}

% Frame: Modelo Pozo Finito
\begin{frame}{5. Modelo de Pozo Finito (Método Numerov)}
\begin{columns}[T]
\column{0.45\textwidth}
\textbf{Definición:}
$$ V(x) = \begin{cases} 0 & |x| < a \\ V_0 & |x| \ge a \end{cases} $$

\textbf{Método Numérico (Numerov):}
Resolvemos la ecuación de Schrödinger independiente del tiempo usando el esquema de diferencias finitas de 4to orden:
$$ \psi_{i+1} = \frac{2(1 - \frac{5}{12}h^2 k_i^2)\psi_i - (1 + \frac{1}{12}h^2 k_{i-1}^2)\psi_{i-1}}{1 + \frac{1}{12}h^2 k_{i+1}^2} $$

\column{0.55\textwidth}
\vspace{0.7cm}
\begin{center}
    \includegraphics[width=0.75\columnwidth]{figures/finite_well_fig1a_Pth.png}
\end{center}
\end{columns}
\end{frame}

% Frame: Resultados Pozo Finito
\begin{frame}{6. Resultados: Pozo Finito}
\begin{columns}[T]
\column{0.55\textwidth}
\vspace{-0.6cm}
\begin{center}
    \includegraphics[width=0.8\columnwidth]{figures/finite_well_fig2_bosons.png}
\end{center}    
\tiny{Bosones dos particulas}
\begin{center}
    \includegraphics[width=0.8\columnwidth]{figures/finite_well_fig3_fermions.png}
\end{center}
\tiny{Fermiones dos particulas}

\column{0.45\textwidth}
\vspace{-1.7cm}
\begin{center}
    \includegraphics[width=0.7\columnwidth]{figures/finite_well_fig4_thermal_2p.eps}
\end{center}
\tiny{(a) Bosones (atracción), (b) Fermiones (repulsión), (c) Spin.}
\end{columns}
\end{frame}

\begin{frame}{6.Resultados: Pozo finito}
\begin{center}
    \includegraphics[width=0.85\columnwidth]{figures/finite_well_fig5_spin.png}
\end{center}    
\end{frame}

% Frame: Extensión Potenciales
\begin{frame}{7. Extensión: Potenciales Armónico y en V}
\begin{columns}[T]
\column{0.25\textwidth}
\textbf{Modelo Armónico Finito:}
$$ V(x) = \frac{1}{2}kx^2 $$

\textbf{Modelo en V Finito:}
$$ V(x) = V_0 \frac{|x|}{a} $$


\column{0.75\textwidth}
\vspace{-0.5cm}
\begin{center}
    \includegraphics[width=0.45\columnwidth]{figures/harmonic_pot.eps}
    \includegraphics[width=0.45\columnwidth]{figures/v_pot.eps}
\end{center}
\end{columns}
\end{frame}

\begin{frame}{8. Comparación de Estadística en Potenciales}

\vspace{-0.2cm}
\begin{columns}[T]
\column{0.35\textwidth}
\centering \textbf{Pozo Finito}
\includegraphics[width=\columnwidth]{figures/finite_well_fig1a_Pth.png}
\tiny{T_{th}}

\column{0.35\textwidth}
\centering \textbf{Armónico}
\includegraphics[width=\columnwidth]{figures/harmonic_fig1a_Pth.png}
\tiny{T_{th}}

\column{0.35\textwidth}
\centering \textbf{Potencial V}
% Placeholder si no existe, o usar harmonic fermiones
\includegraphics[width=\columnwidth]{figures/harmonic_fig1a_Pth.png}
\tiny{T_{th}}
\end{columns}
\end{frame}


% Frame: Resultados Comparados
\begin{frame}{8. Comparación de Estadística en Potenciales}

\vspace{-0.2cm}
\begin{columns}[T]
\column{0.35\textwidth}
\centering \textbf{Pozo Finito}
\includegraphics[width=\columnwidth]{figures/finite_well_fig2_bosons.png}
\tiny{Bosones}

\column{0.35\textwidth}
\centering \textbf{Armónico}
\includegraphics[width=\columnwidth]{figures/harmonic_fig2_bosons.png}
\tiny{Bosones}

\column{0.35\textwidth}
\centering \textbf{Potencial V}
% Placeholder si no existe, o usar harmonic fermiones
\includegraphics[width=\columnwidth]{figures/harmonic_fig2_bosons.png}
\tiny{Bosones}
\end{columns}
\end{frame}


% Frame: Resultados Comparados
\begin{frame}{8. Comparación de Estadística en Potenciales}

\vspace{-0.2cm}
\begin{columns}[T]
\column{0.35\textwidth}
\centering \textbf{Pozo Finito}
\includegraphics[width=\columnwidth]{figures/finite_well_fig3_fermions.png}
\tiny{Fermiones}
\column{0.35\textwidth}
\centering \textbf{Armónico}
\includegraphics[width=\columnwidth]{figures/harmonic_fig3_fermions.png}
\tiny{Fermiones}
\column{0.35\textwidth}
\centering \textbf{Potencial V}
% Placeholder si no existe, o usar harmonic fermiones
\includegraphics[width=\columnwidth]{figures/harmonic_fig3_fermions.png}
\tiny{Fermiones}
\end{columns}
\end{frame}


\begin{frame}{9.Convergencia}
\begin{center}
    \includegraphics[width=0.35\columnwidth]{figures/convergence_t_vs_N.png}
    \includegraphics[width=0.35\columnwidth]{figures/convergence_fit_t_vs_N.png}
    \includegraphics[width=0.35\columnwidth]{figures/convergence_t_vs_N_clean.png}
\end{center}
\end{frame}

\begin{frame}{10.Densidad de Estados}
\begin{center}
    \includegraphics[width=0.35\columnwidth]{figures/dos_comparison_multiple_widths.png}
    \includegraphics[width=0.35\columnwidth]{figures/dos_convergence_comparison.png}
    \includegraphics[width=0.35\columnwidth]{figures/dos_cumulative_N.png}
\end{center}
\end{frame}

%===============================================================================
% CONCLUSIONES
%===============================================================================
\section{Conclusiones}
\begin{frame}{Conclusiones}
\begin{itemize}
\item \textbf{Efecto Térmico Universal:} En todos los potenciales, $T$ tiende a uniformizar la densidad $P_{th}$, pero las condiciones de frontera (duras o suaves) imponen perfiles específicos en los bordes.
\item \textbf{Robustez Estadística:} La naturaleza de Bosones (atracción efectiva) y Fermiones (repulsión) domina sobre la geometría del potencial.
\item \textbf{Espín:} Es el mecanismo fundamental que permite a los fermiones (como electrones) compartir espacio mediante el estado Singlete.
\item \textbf{Métodos Numéricos:}
    \begin{itemize}
        \item \textit{Numerov:} Alta precisión para estados 1D continuos.
    \end{itemize}
\end{itemize}

\vspace{0.3cm}

\end{frame}

%===============================================================================
% DIAPOSITIVA FINAL
%===============================================================================
\begin{frame}[plain]
\begin{center}
\Huge{\textcolor{udea}{\textbf{Gracias}}}

\vspace{1cm}
\normalsize
\textbf{Juan Montoya \& Pablo Sánchez}\\
Física Estadística
\end{center}
\end{frame}

\end{document}
